\starttext
\startplacetable[location={split},title={Begriffe}]
\setupTABLE[c][1][background=color,backgroundcolor=gray]
\setupTABLE[r][1][background=color,backgroundcolor=gray]
\bTABLE[split=yes,option=stretch,frame=on]
\bTR
\bTD Begriff \eTD  \bTD Definition \eTD
\eTR

\bTR
\bTD Zufallsexperiment: \eTD
\bTD
Ein Vorgang, welcher unter den selben Bedingungen beliebig oft wiederholt werden kann, wobei dessen Ausgang nicht vorhersehbar ist. Beispiel: Würfel, Münzwurf.
\eTD
\eTR

\bTR
\bTD Ergebnis: \eTD
\bTD
Ein möglicher Ausgang eines Zufallsexperiment. Beispiel: Alle Zahlen zwischen Eins und Sechs (${1;2;3;4;5;6}$) beim würfeln eines 6-seitiegen Würfel.
\eTD
\eTR

\bTR
\bTD Ereignis: \eTD
\bTD
Eine {\em Teilmenge} der Ereignismenge.
\eTD
\eTR

\bTR
\bTD Wahrscheinlichkei ($P$): \eTD
\bTD Eine Maß, wie wahrscheinlich das Eintreten eines Ereignisses ist. Angegeben mit einer Zahl von Null bis Eins ($\[0;1\]$)
\eTD
\eTR

\bTR
\bTD Faires Spiel: \eTD
\bTD 
Ein Spiel, bei dem der Erwartungswert für den Gewinn gleich Null ist:

\startformula
\text{Gewinn} = \text{Auszahlung} - \text{Einsatz}
\stopformula

Beispiel: Beim Würfeln soll das doppelte der Augenzahl in Euro ausgezahlt werden. Die Ergebnismenge beinhaltete die möglichen Auszahlungen:

\startformula
{2;4;6;8;10;12}
\stopformula

Somit entspricht der Erwartungswert:

\startformula
\mu = \frac{1}{6}\cdot (2+4+6+8+10+12) = \frac{42}{6}=7
\stopformula

Somit muss der Einsatz $7$ betragen, damit es sich um ein \quotation{Faires Spiel} handelt
\eTD
\eTR

\bTR
\bTD Laplace - Experiment: \eTD
\bTD
Ein Versuch, bei dem alle Ereignisse gleich wahrscheinlich sind. Beispiel: Würfeln
\eTD
\eTR

\bTR
\bTD Unabhägigkeit \eTD
\bTD
Wenn das Eintreten des einen Ereignisses die Wahrscheinlichkeit des anderen {\bf nicht} beeinflusst.
\eTD
\eTR

\bTR
\bTD Wahrscheinlichkeitsverteilung: \eTD
\bTD 
Wenn jedem Ereignis eines Zufallsexperiment ein Zahlenwert zugeordnet wird, spriucht man von einer Zufallsgröße. Die Wahrscheinlichkeitsverteilung einer Zufallsgröße $X$, ist eine Tabelle, bei der jedem Wert von $k$ von $X$ die Wahrscheinlichkeit $P(X=k)$ zugeordnet ist.

\framed[aligne=middle,frame=off,offset=none]{
\starttabulate[|m|c|c|c|c|c|]
\NC k \VL 0 \NC 1 \NC 2 \NC 3 \NC 4 \NC\NR
\HL
\NC P(X=k) \VL 0,16 \NC 0,32 \NC 0,29 \NC 0,15 \NC 0,05 \NC \NR
\stoptabulate
}
\eTD
\eTR

\bTR
\bTD Varianz ($V$) \eTD
\bTD
ist eine Maß für die Streung einer Zufallsgröße.

\startformula
V=(x_1 - \mu)^2 \cdot P(X=x_1) + \dots + (x_n - \mu)^2 \cdot P(X=x_n)
\stopformula

\eTD
\eTR
\eTABLE
\stopplacetable


\stoptext
